% ============================================================================
%  Section VII — Discussion and Limitations  (IOP MST)
%  Passive voice / third person throughout. Interpretation of the Section VI
%  results, practical guidance, and limitations. Requires: none beyond base.
% ============================================================================

\section{Discussion and Limitations}
\label{sec:discussion}

\subsection{Discussion}

\subsubsection{Evaluation methodology}
The optimism gap of Section~\ref{subsec:gap} is the central methodological
finding: for low-rate smartphone-based machine diagnosis, accuracies obtained
under random splitting overstate deployable performance by tens of accuracy
points. The effect mirrors that documented for industrial-grade vibration
data~\cite{wheat2024impact, hendriks2022towards} and the broader
machine-learning leakage literature~\cite{kapoor2023leakage, lones2024leakage},
and it has been observed independently in wearable
sensing~\cite{gholamiangonabadi2020cv}.

The decomposition of Section~\ref{subsec:decomp} sharpens the usual account.
Window overlap, which is the mechanism most often named and most often guarded
against, accounts for none of the gap here: a random split of strictly
non-overlapping windows is still perfectly accurate, as is a split that holds out
a contiguous temporal block from every recording. The whole of the $46.7$-point
gap is attributable to recording identity. The practical implication is
uncomfortable for common practice, because enforcing non-overlapping windows---a
frequent and reasonable-looking precaution---would provide no protection at all on
data of this kind. Where each operating condition is captured in one continuous
acquisition, the recording is the only defensible unit of separation.

A consequence for benchmark practice follows. Near-perfect accuracies frequently
reported for this class of data should be read as measurements of recording
identifiability rather than of diagnostic capability, a concern raised
independently for deep bearing-diagnosis benchmarks~\cite{wang2026benchmark}. The
recording-wise and cross-condition figures obtained here provide a more
deployment-honest reference against which future methods can be compared.

\subsubsection{Sampling rate, feature invariance and what the sweep measures}
Decimation to $25$~Hz raises leakage-free accuracy, and the resource-frugal
configuration built on it remains the better operating point for a node. The
reason, however, is not the one a rate sweep invites. Section~\ref{subsec:antialias}
shows that band-limiting to the same $12.5$~Hz while retaining the original
sampling rate \emph{reduces} accuracy, that the whole improvement appears when the
band-limited signal is represented by a quarter as many samples per window, and
that disabling the anti-alias filter is severely damaging. The shaft rotational
frequencies of this machine, $13.1$ to $24.9$~Hz, lie above the $12.5$~Hz Nyquist
frequency of the decimated stream and are therefore absent from the configuration
that performs best.

The operative effect is that the time-domain feature set is not invariant to the
sampling rate. Zero-crossing rate, amplitude entropy over a fixed histogram and
the higher-order moments all depend on the number of samples in the window and on
how strongly those samples are correlated; a band-limited signal retained at
$100$~Hz is heavily oversampled, and its descriptors are estimated from far fewer
effectively independent values than the sample count implies. Decimation removes
that redundancy, and the gain follows.

Two conclusions are worth separating. The practical one is unchanged and useful:
an ultra-low-rate node is not handicapped for this task, and aggressive decimation
reduces telemetry and node energy without sacrificing accuracy---provided the
anti-alias filter is applied, since omitting it costs $24.9$ points. The
methodological one is a caution that applies well beyond this dataset. A sampling
rate sweep performed with a fixed time-domain feature set does not measure the
information content of the band; it measures the band and the sample-count
sensitivity of the features together. Reports that lower rates suffice, this one
included, should be read with that confound in mind, and studies of downsampling
in vibration analysis~\cite{he2021sampling} suggest the interaction deserves
explicit treatment.

\subsubsection{Model complexity and edge feasibility}
Under leakage-free evaluation the lower-capacity models generalised at least as
well as the Random Forest while occupying a kilobyte-scale footprint, so on-node
inference is practical on inexpensive microcontrollers. Combined with the link
budget, this yields per-link deployment guidance: Bluetooth Low Energy, Zigbee,
and NB-IoT nodes can stream raw or feature data to a gateway, whereas LoRaWAN
nodes, constrained by the duty cycle to the decision stream, must perform
classification locally. The alignment between the models that deploy most easily
and those that generalise best is convenient for practical systems.

\subsubsection{Sensor-level guidance}
The channel analysis of Section~\ref{sec:charact} carries a concrete
recommendation for smartphone-based acquisition: the raw acceleration channels
should be preferred over the gravity-compensated channels. The on-device gravity
compensation---valuable for motion-tracking applications---preserves the vibration
power but reshapes its low-frequency structure, and the raw channels carry
measurably more fault information (a mutual-information ratio of about $1.4$ that
survives the removal of static offsets), so they are the better choice for
condition monitoring.

\subsubsection{Scope of applicability}
Taken together, the results delineate where low-cost, low-rate MEMS sensing is
appropriate. Coarse health states, such as a healthy machine, a broken rotor bar,
or a supply-voltage imbalance, are separable at $100$~Hz and below; the staging of
incipient bearing degradation, whose signatures lie in the high-frequency band, is
not, and remains the province of higher-rate, industrial-grade instrumentation.

\subsection{Limitations}
\label{subsec:limits}
Several limitations bound the scope of these findings. The dataset comprises a
single $1.1$~kW machine with laboratory-induced, discrete-severity faults, and the
two lubrication states were graded qualitatively; the reported accuracies are
therefore a deployment-honest baseline rather than a claim of field performance.
Crucially, each fault class is represented by a \emph{single physical specimen}
(one modified bearing, one drilled rotor, one imbalance configuration) recorded
across speeds and loads. The recording-wise protocol therefore removes
window-overlap and exact operating-point leakage, but not specimen identity: a
model may still key on the idiosyncratic signature of that particular specimen and
its mounting rather than on a fault signature that would transfer to a different
specimen of the same class. The evaluation is thus leakage-free only in this
restricted sense, and cross-specimen generalisation is not established here. The
small number of recordings ($36$) also produces wide fold-to-fold variability;
although five-seed repetition stabilises the central estimates and the reported
expanded uncertainty quantifies their reproducibility, it captures only resampling
variance and not the dominant specimen-, machine-, and mounting-to-mounting
components, so it should not be read as a full measurement-uncertainty budget. The
deployment quantities are likewise indicative rather than measured: the model
footprint is a serialised (software) size, and energy and latency are inferred
from data volume and window length rather than measured on hardware, with a
cycle-accurate microcontroller characterisation left to future work. A further limitation concerns
reproducibility in practice: the implementation itself is not distributed, so the
results must be reimplemented from the description rather than re-executed. The
feature definitions, estimator settings, resampling and robustness transforms,
and random seeds are therefore reported in full in Section~\ref{sec:method},
but exact numerical agreement with an independent reimplementation cannot be
guaranteed. Finally,
generalisation across machines of different ratings and manufacturers was not
assessed and constitutes the most important direction for future validation.
